\section{With the sources unknown}
\label{sec:blind}

\begin{figure*}[t]
  \centering
  \includegraphics[width=\textwidth]{fig2_diagnostics.pdf}
  \caption{What a practitioner can compute without ground truth, against what the benchmark knows. (a) Spearman correlation between an observational diagnostic of the fitted model and its true held-out recovery, pooled over the cells of Figure~\ref{fig:axes} and within each stress axis; blue negative, orange positive. (b--d) Three of those diagnostics against true recovery, one point per (axis, condition, seed) cell, coloured by axis: TCL's training segment accuracy, the iVAE evidence lower bound, and the agreement between two TCL fits on the same data.}
  \label{fig:diag}
\end{figure*}

A practitioner has none of Figure~\ref{fig:axes}.
What they have is the fitted model: TCL's segment-classification accuracy, the iVAE's evidence lower bound, the agreement between two fits, the non-Gaussianity of FastICA's components, the explained variance of PCA.
Figure~\ref{fig:diag} correlates each of these with the true held-out recovery across the cells of the benchmark, pooled and within each axis.
No diagnostic keeps a stable relation to recovery.
The accuracy that TCL papers report is anti-diagnostic in training: it is highest where recovery is lowest, because drift is easier to classify than source variance (Figure~\ref{fig:diag}b).
The iVAE's bound tracks recovery on the coupling axis and anti-tracks it on the source-count axis (Figure~\ref{fig:diag}c).
Run-to-run agreement, the best pooled diagnostic, is uninformative inside the nonstationarity axis, where the real nuisance lives (Figure~\ref{fig:diag}d); an iVAE that is consistently wrong is consistent.
The sign of every column of Figure~\ref{fig:diag}a depends on which stress the data happen to carry, so calibrating any of these diagnostics requires the benchmark, and the calibration flips across the stresses that a real recording mixes.
This is the second half of Section~\ref{sec:argument} as a measurement.
